\documentclass[a4paper]{article}
% Fonts: Palatino text + Euler math (Vakil style)
\usepackage[T1]{fontenc}
\usepackage{newpxtext}
\usepackage[small]{eulervm}
\linespread{1.05}

% Layout
\usepackage[a4paper,margin=2cm]{geometry}
\usepackage{microtype}
\setlength{\parindent}{0pt}

% Math
\usepackage{amsthm,amsmath,amssymb}

% Utilities
\usepackage{graphicx}
\usepackage{tcolorbox}
\usepackage[colorlinks,allcolors=black]{hyperref}

% Header/footer
\usepackage{fancyhdr}
\pagestyle{fancy}
\fancyhf{}
\renewcommand{\headrulewidth}{0pt}
\fancyfoot[R]{\thepage}

\renewcommand{\thesubsection}{\alph{subsection}}

\DeclareMathOperator{\Res}{Res}
\DeclareMathOperator{\Ind}{Ind}
\DeclareMathOperator{\Log}{Log}
\newcommand{\C}{\mathbb{C}}
\newcommand{\Z}{\mathbb{Z}}
\newcommand{\N}{\mathbb{N}}
\newcommand{\Rb}{\mathbb{R}}

\newtcolorbox{prob}{colback=gray!10, colframe=black, boxrule=0.5pt, arc=2mm}

\begin{document}

\hrule\medskip
\begin{minipage}{0.295\textwidth}\raggedright\footnotesize
Ali Azam
\end{minipage}%
\begin{minipage}{0.4\textwidth}\centering
\large Math 113: Complex Analysis\\
\normalsize HW \#8
\end{minipage}%
\begin{minipage}{0.295\textwidth}\raggedleft
\today
\end{minipage}
\medskip\hrule\bigskip

\section*{Ex 9.1.1}

\begin{prob}
Prove Proposition 9.2. (Hint: Either use Proposition 9.3 or use the observation
that the formula (9.4) holds for all circles $\gamma_\tau$ with $\tau$
sufficiently small.)

\medskip
\emph{Proposition 9.2.} Suppose that $f$ has an isolated singularity at $a$ with
associated Laurent series
\[
f(z)=\sum_{n=0}^{\infty}a_n(z-a)^n+\sum_{j=1}^{\infty}\frac{b_j}{(z-a)^j}.
\]
\begin{itemize}
\item[(a)] $a$ is a removable singularity if and only if all the $b_j$ are zero.
\item[(b)] $a$ is a pole of order $m\ge 1$ if and only if $b_m\ne 0$ and $b_j=0$
for $j>m$.
\item[(c)] $a$ is an essential singularity if and only if $b_j\ne 0$ for
infinitely many $j$.
\end{itemize}
\end{prob}

\textit{Pf.} We have that the Laurent series above is valid on a punctured disk, $0<|z-a|<r$. We use Thm~9.1's (Laurent's Theorem) uniqueness statement to say that the coefficients $a_n,b_j$ are determined by $f$.

\medskip

\begin{enumerate}
    \item[\textbf{(a)}] ($\Leftarrow$) Let $\forall b_j = 0$, then $f(z) = \sum_{n\ge0}a_n(z-a)^n$ on the punctured disk. Using Thm 1.34, this power series defines a holomorphic function on the full disk $D_r(a)$, therefore setting $f(a) = a_0$ makes $f$ holomorphic at $a$. Therefore, by Def~5.12, the singularity is removable.

    \medskip
    ($\Rightarrow$) If $a$ is removable, let $g$ be the holomorphic extension of $f$ to $D_r(a)$. By Cor~4.14 (Taylor's Theorem), we have on $D_r(a)$,
    \[
    g(z) = \sum_{n\ge0}c_n(z-a)^n
    \]  
    which is the Laurent series for $f$ on $0<|z-a|<r$ whose negative coefficients all vanish. Therefore, using uniquness in Thm~9.1 we must have $\forall b_j = 0$ and $a_n = c_n$. \qed

    \item[\textbf{(b)}] ($\Leftarrow$) Let $b_m \ne 0$ and $b_j = 0$ for $j > m$. Then, we have,
    \[
    f(z) - \sum_{j=1}^m \frac{b_j}{(z-a)^j} = \sum_{n\ge0}a_n(z-a)^n
    \]
    which is known to be removable by part (a). By Def~5.14, we have exactly that the $a$ is pole of order $m$.

    \medskip
    ($\Rightarrow$) Let $a$ be a pole of order $m$. By Def~5.14, we have $\exists c_1, \dots, c_m$ with $c_m \ne 0$ s.t.
    \[
    r(z) = f(z) - \sum_{j=1}^m c_j(z-a)^{-j}
    \]
    has a removable singularity at $a$. Using part (a), the Laurent series of $r$ has no principal part, therefore,
    \[
    f(z) = \sum_{n\ge0} a_n(z_a)^n + \sum_{j=1}^m \frac{c_j}{(z-a)^j}
    \]
    By Thm 9.1's uniqueness, $\forall j, 1 \le j \le m: b_j = c_j$ and $\forall j > m: b_j = 0$. I.e., we have $b_m = c_m \ne 0$. \qed

    \item[\textbf{(c)}] ($\Rightarrow$) By Def~5.20, $a$ is essential when it is neither removable nor a pole. By part (a), if $a$ is removable then $\forall b_j = 0$, and by part (b) if $a$ is a pole then it has the largest index $m$ with $b_m \ne 0$ s.t. $\forall j > m: b_j =0$. In either case, only finitely many $b_j$ are nonzero.

    \medskip
    ($\Leftarrow$) Ff only finitely many $b_j$ aare nonzero, then either all vanish (removabel) or there is a largest $m$ with $b_m \ne 0$ (pole of order $m$). 

    \medskip
    Therefore, $a$ is essential iff infinitely many $b_j$ are nonzero. \qed
\end{enumerate}

\newpage
\section*{Ex 9.2.2}

\begin{prob}
Suppose that $f$ has an isolated singularity at $a$. Show that if
$\lim_{z\to a}(z-a)f(z)=L\ne 0$, then $f$ has a simple pole at $a$ with residue
$L$. What happens if $L=0$?
\end{prob}

\textit{Pf.} On the punctured disk where $f$, we consider $\varphi(z)=(z-a)f(z)$. We can see that $\lim_{z\to a}\varphi(z)=L$ exists and is finite (since $\varphi$ is bounded near $a$). This fulfills Prop~5.13's boundless criterion (d), since the limit exists, which gives us (a) i.e. the singularity of $\varphi$ at $a$ is removable.

\medskip
We can then extend $\varphi$ holomorphically s.t. $\varphi(a)=L$, and can write its Taylor series as,
\[
\varphi(z)=\sum_{n\ge 0}c_n(z-a)^n
\]
Therefore, $c_0=\varphi(a)=L$. On the punctured disk, we have,
\[
f(z)=\frac{\varphi(z)}{z-a}=\frac{L}{z-a}+\sum_{n\ge 1}c_n(z-a)^{n-1}
\]
where the second term is a power series and holomorphic at $a$. Therefore, the Laurent series of $f$ at $a$ has $b_1=L\ne 0$ and $\forall j \ge 2: b_j=0$. 

\medskip
Using Prop~9.2 (b), proven above in Ex~9.1.1, we have $a$ is a pole of order $1$, i.e. a simple pole. By Def~9.6 we also have that the residue is $\Res(f;a)=b_1=L$. \qed

\medskip
\underline{In the case $L=0$}: using the same procedure gives $\varphi$ holomorphic at $a$ with $\varphi(a)=0$. Therefore, we have $\varphi(z)=\sum_{n\ge 1}c_n(z-a)^n$ and that,
\[
f(z)=\frac{\varphi(z)}{z-a}=\sum_{n\ge 1}c_n(z-a)^{n-1}=\sum_{m\ge 0}c_{m+1}(z-a)^m
\]
is a power series. Therefore $f$ has a removable singularity at $a$, i.e. $a$ is not a pole. So when $L=0$ the only possibility is that the
singularity is removable.

\newpage
\section*{Ex 9.2.3}

\begin{prob}
Prove the Simple Pole Lemma on page 138.

\medskip
\textbf{Lemma 9.16 (Simple Pole Lemma).} Suppose that $g$ and $h$ are holomorphic
at $a$, that $g(a)\ne 0$ and that $h$ has a simple zero at $a$. Then $f=g/h$ has
a simple pole at $a$ and
\[
\Res(f;a)=\frac{g(a)}{h'(a)}.
\]
\end{prob}

\textit{Pf.} Since $h$ has a simple zero (a zero of order $1$) at $a$, we have from Lemm~5.4 that the function $h_1$ is holomorphic at $a$ with $h_1(a)\ne 0$ and that,
\[
h(z)=(z-a)\,h_1(z).
\]

We can differentiate to get, 
\[
h'(z)=h_1(z)+(z-a)h_1'(z)
\]
Therefore, $h'(a)=h_1(a)\ne 0$. Since $h_1(a)\ne 0$ and $g(a)\ne 0$, we have a punctured disk on which $h_1$, and hence $h$, are nonzero. Therefore, $f=g/h$ is holomorphic there and has an isolated singularity
at $a$. Then,
\[
\lim_{z\to a}(z-a)f(z)
=\lim_{z\to a}\frac{(z-a)g(z)}{(z-a)h_1(z)}
=\lim_{z\to a}\frac{g(z)}{h_1(z)}
=\frac{g(a)}{h_1(a)}=\frac{g(a)}{h'(a)}\ne 0
\]
From Ex~9.2.2 above we have that $f$ has a simple pole at $a$ with residue equal to this limit, i.e. $\displaystyle \Res(f;a)=\frac{g(a)}{h'(a)}$. \qed

\newpage
\section*{Ex 9.2.5}

\begin{prob}
Suppose that $f$ has a pole of order $m$ at $z_0$ and that $k$ has a zero of
order $n$ at $z_0$. Show that $\Res(f'/f;z_0)=-m$ and that $\Res(k'/k;z_0)=n$.
(Compare with Ex 5.2.7.)
\end{prob}

\textit{Pf.} Since $k$ has a zero of order $n$ at $z_0$, using Lemm~5.4 gives a function $K$ holomorphic at $z_0$ with $K(z_0)\ne 0$ and $k(z)=(z-z_0)^nK(z)$. Then we can differentiate to get,
\[
k'(z) = n(z-z_0)^{n-1}K(z)+(z-z_0)^nK'(z)
\]
which gives,
\[
\frac{k'(z)}{k(z)} = \frac{n}{z-z_0}+\frac{K'(z)}{K(z)}
\]
Since $K(z_0)\ne 0$, the term $K'/K$ is holomorphic near $z_0$. Hence $k'/k$ has a simple pole at $z_0$ (as in Ex~5.2.7). We can use $\lim_{z\to z_0}(z-z_0)\,k'/k=n\ne 0$ with Ex~9.2.2 to obtain,
\[
\Res\!\left(\frac{k'}{k};z_0\right)=n.
\]

\medskip
Since $f$ has a pole of order $m$ at $z_0$, Cor~5.18
gives a function $F$ holomorphic at $z_0$ with $F(z_0)\ne 0$ and
$f(z)=F(z)/(z-z_0)^m$. Therefore we have,
\begin{align*}
f'(z)&=\frac{F'(z)}{(z-z_0)^m}-\frac{mF(z)}{(z-z_0)^{m+1}} \\
\frac{f'(z)}{f(z)}&=\frac{F'(z)}{F(z)}-\frac{m}{z-z_0}.
\end{align*}
Since $F(z_0)\ne 0$, $F'/F$ is holomorphic near $z_0$. Hence $f'/f$ has a simple pole at $z_0$, and its residue is the coefficient of $(z-z_0)^{-1}$, i.e.
$\displaystyle \Res\!\left(\frac{f'}{f};z_0\right)=-m$. \qed

\newpage
\section*{Ex 9.3.5}

\begin{prob}
Show that
\[
I=\int_{0}^{2\pi}\cos(\theta)^{2n}\,d\theta=\frac{\pi(2n)!}{2^{2n-1}(n!)^2}.
\]
\end{prob}

\textit{Pf.} Like in Section~9.3.1, parametrize $z=e^{i\theta}$, s.t. $\theta\mapsto z$ traverses the
positively oriented unit circle $\gamma$, with $dz=iz\,d\theta$ (hence
$d\theta=dz/(iz)$) and $\cos\theta=\tfrac12(z+z^{-1})$. Therefore we have,
\[
I=\oint_{\gamma}\left(\frac{z+z^{-1}}{2}\right)^{2n}\frac{dz}{iz}
=\frac{1}{2^{2n}i}\oint_{\gamma}\frac{(z+z^{-1})^{2n}}{z}\,dz
\]
Using the binomial theorem, we have,
\[
(z+z^{-1})^{2n}=\sum_{k=0}^{2n}\binom{2n}{k}z^{2k-2n}
\]
therefore,
\[
\frac{(z+z^{-1})^{2n}}{z}=\sum_{k=0}^{2n}\binom{2n}{k}z^{2k-2n-1}
\]
We have that the only singularity of the integrand is at $0$, and its residue is the coefficient of $z^{-1}$, given by $2k-2n-1=-1 \Rightarrow k=n$:
\[
\binom{2n}{n}=\frac{(2n)!}{(n!)^2}
\]
We have that $\gamma$ is a positively-oriented Jordan path and the integrand is holomorphic inside (besides at $0$), we have using Cor~9.13 (the Residue Theorem) that
$\oint_{\gamma}=2\pi i\binom{2n}{n}$. Therefore,
\[
I=\frac{1}{2^{2n}i} 2\pi i\frac{(2n)!}{(n!)^2}
=\frac{2\pi}{2^{2n}}\frac{(2n)!}{(n!)^2}
=\frac{\pi(2n)!}{2^{2n-1}(n!)^2}
\]
as required. \qed

\newpage
\section*{Ex 10.1.1}

\begin{prob}
Suppose that $f$ is holomorphic on the punctured complex plane
$D^0=\C\setminus\{0\}$. Suppose there is a sequence of polynomials $(p_n)$
converging to $f$ uniformly on compact subsets of $D^0$. Show that $f$ has a
removable singularity at $0$.
\end{prob}

\textit{Pf.} In $D^0=\C\setminus\{0\}$, deleted neighborhood of $0$, we have that $0$ is an isolated singularity of $f$. Using Thm~9.1 (Laurent's Theorem), for $0 < |z| < \infty$, $f$ has a Laurent series,
\[
f(z)=\sum_{n\ge 0}a_nz^n+\sum_{j\ge 1}\frac{b_j}{z^j}
\]
We can find its coefficients using equation~(9.4), $\forall \rho>0$ and $\forall j\ge 1$,
\[
b_j=\frac{1}{2\pi i}\oint_{\gamma_\rho}f(w)\,w^{j-1}\,dw
\]
where $\gamma_\rho$ is the positively-oriented circle $|w|=\rho$.

\medskip
We can consider $\rho>0$, s.t. the circle $\gamma_\rho^{*}$ is a compact subset of $D^0$. Therefore, $p_n\to f$ uniformly there. Since $w^{j-1}$ is bound on $\gamma_\rho^{*}$
we have $p_n(w)w^{j-1}\to f(w)w^{j-1}$ uniformly on $\gamma_\rho^{*}$. 

\medskip
Using Prop~3.26(b) we can pass to the limit inside the integral s.t.
\[
b_j=\lim_{n\to\infty}\frac{1}{2\pi i}\oint_{\gamma_\rho}p_n(w)\,w^{j-1}\,dw
\]
We can consider $j\ge 1$, s.t. the integrand $p_n(w)w^{j-1}$ is a polynomial. Using Lemm~4.5 (polynomial has an antiderivative, and $\gamma_\rho$ is closed), its integral over $\gamma_\rho$ is $0$. Therefore, $\forall j \ge 1: b_j=0$. 

\medskip
Using Prop~9.2(a) (proved in Ex~9.1.1 above), the singularity of $f$ at $0$ is removable. \qed

\newpage
\section*{Ex 9.3.10}

\begin{prob}
Evaluate
\[
\int_{0}^{2\pi}\frac{\cos(\theta)}{3+\cos(\theta)}\,d\theta.
\]
\end{prob}

\textit{Pf.} We can rewrite the integral as $\dfrac{\cos\theta}{3+\cos\theta}=1-\dfrac{3}{3+\cos\theta}$, s.t.
\begin{align*}
J: &=\int_{0}^{2\pi}\frac{d\theta}{3+\cos\theta} \\
\\
I &=\int_{0}^{2\pi}\frac{\cos\theta}{3+\cos\theta}\,d\theta
=\int_{0}^{2\pi}d\theta-3\int_{0}^{2\pi}\frac{d\theta}{3+\cos\theta}
=2\pi-3J
\end{align*}
For $J$ we parametrize $z=e^{i\theta}$, s.t. $d\theta=dz/(iz)$, $\cos\theta=\tfrac12(z+z^{-1})$
and $3+\cos\theta=\dfrac{z^2+6z+1}{2z}$. Then we can consider it over the positively oriented unit circle $\gamma$ so that we have,
\[
J=\oint_{\gamma}\frac{2z}{z^2+6z+1}\frac{dz}{iz}
=\frac{2}{i}\oint_{\gamma}\frac{dz}{z^2+6z+1}
\]
We have that the denominator has roots $z_{\pm}=-3\pm 2\sqrt2$. Since $z_+z_-=1$ and
$|z_-|=3+2\sqrt2>1$, only $z_+=-3+2\sqrt2$ lies inside $\gamma$, and it is a
simple pole. Using Ex~9.2.3 (done above) and Lemm~9.16, with denominator $z^2+6z+1$ of derivative $2z+6$, we have,
\[
\Res\!\left(\frac{1}{z^2+6z+1};z_+\right)=\frac{1}{2z_++6}
=\frac{1}{2(-3+2\sqrt2)+6}=\frac{1}{4\sqrt2}
\]
Using Cor~9.13 (the Residue Theorem), we have,
\[
J=\frac{2}{i} 2\pi i\frac{1}{4\sqrt2}=\frac{\pi}{\sqrt2}
\]
Therefore, we have the integral evaluation as,
\[
I=2\pi-3\cdot\frac{\pi}{\sqrt2}=2\pi-\frac{3\sqrt2}{2}\pi
=\frac{(4-3\sqrt2)\pi}{2}
\]

\newpage
\section*{Ex 9.3.11}

\begin{prob}
Evaluate
\[
\int_{0}^{\infty}\frac{\cos(ax)}{(x^2+b^2)^2}\,dx
\qquad\text{where }a>0\text{ and }b>0.
\]
\end{prob}

\textit{Pf.} We have that the integrand is even, therefore $\int_0^\infty=\tfrac12\int_{-\infty}^{\infty}$. Let
\begin{align*}
f(z)&=\frac{1}{(z^2+b^2)^2}=\frac{1}{(z-ib)^2(z+ib)^2}
\\ 
g(z)&=f(z)e^{iaz}
\end{align*}
We have $f=p/q$ with $\deg p=0$ and $\deg q=4$, therefore $\deg p+2\le\deg q$, and $q$ has
no real zeros. Since $a>0$, we can use Prop~9.21 to get,
\[
\int_{-\infty}^{\infty}\frac{e^{iax}}{(x^2+b^2)^2}\,dx
=2\pi i\sum_{\operatorname{Im}z>0}\Res(g;z)=2\pi i\,\Res(g;ib)
\]
where the only pole in the upper half-plane is the double pole at $ib$. Using Lemm~9.17, we have,
\[
\Res(g;ib)=\lim_{z\to ib}\frac{d}{dz}\Big[(z-ib)^2g(z)\Big]
=\frac{d}{dz}\left[\frac{e^{iaz}}{(z+ib)^2}\right]_{z=ib}
\]
We have,
\[
\frac{d}{dz}\Big[e^{iaz}(z+ib)^{-2}\Big]
=e^{iaz}(z+ib)^{-3}\big[\,ia(z+ib)-2\,\big]
\]
At $z=ib$ we have $z+ib=2ib$, $e^{iaz}=e^{-ab}$, $ia(2ib)-2=-2(ab+1)$, and
$(2ib)^{-3}=\dfrac{i}{8b^3}$. Therefore,
\[
\Res(g;ib)=e^{-ab}\cdot\frac{i}{8b^3}\cdot(-2)(ab+1)
=-\,\frac{i(ab+1)e^{-ab}}{4b^3}
\]
and we get,
\[
\int_{-\infty}^{\infty}\frac{e^{iax}}{(x^2+b^2)^2}\,dx
=2\pi i\cdot\left(-\frac{i(ab+1)e^{-ab}}{4b^3}\right)
=\frac{\pi(1+ab)e^{-ab}}{2b^3}.
\]
Using Rmk~9.23, we take the real parts and replace $e^{iax}$ by $\cos(ax)$. The value above is already real, therefore,
\[
\int_{-\infty}^{\infty}\frac{\cos(ax)}{(x^2+b^2)^2}\,dx
=\frac{\pi(1+ab)e^{-ab}}{2b^3}
\]
We can halve this to obtain the integral evaluation as,
\[
\int_{0}^{\infty}\frac{\cos(ax)}{(x^2+b^2)^2}\,dx
=\frac{\pi(1+ab)e^{-ab}}{4b^3}
\]

\newpage
\section*{Ex 9.3.12}

\begin{prob}
Evaluate $\displaystyle\int_{0}^{\infty}\frac{1}{x^n+1}\,dx$ where $n\ge 2$.
(Note that we cannot use the old reliable contour $\Gamma_R$ from Proposition
9.21 since the integrand has a singularity at $x=-1$ when $n$ is odd. Instead, I
suggest you use the contour which is the positively oriented boundary of the
circular sector $S=\{re^{i\theta}:0\le\theta\le\tfrac{2\pi}{n}\text{ and }
0\le r\le\alpha\}$.)
\end{prob}

\textit{Pf.} Let $f(z)=\dfrac{1}{z^n+1}$ and let $\Gamma$ be the positively oriented boundary
of the sector $S$ consisting of the segment $[0,\alpha]$ on the positive real axis, the arc
$C_\alpha$ given by $z=\alpha e^{i\theta}$ for $0\le\theta\le 2\pi/n$, and the
segment back to the origin along the ray $\arg z=2\pi/n$.

\medskip
We have that the zeros of $z^n+1$ are $z=e^{i\pi(2k+1)/n}$. The only zero with argument in
$(0,2\pi/n)$ is $z_0=e^{i\pi/n}$ (the next zero, $e^{i3\pi/n}$, has argument $3\pi/n>2\pi/n$). Therefore $z_0$ is the unique pole inside the sector, and it is
simple. Using Ex~9.2.3 (done above) and Lemm~9.16 (Simple Pole Lemma), with denominator
$z^n+1$ of derivative $nz^{n-1}$, and using $z_0^n=-1$, we have,
\[
\Res(f;z_0)=\frac{1}{nz_0^{n-1}}=\frac{z_0}{nz_0^{n}}=\frac{z_0}{-n}
=-\frac{e^{i\pi/n}}{n}
\]
Using Cor~9.13 (the Residue Theorem), we have,
\[
\oint_{\Gamma}f(z)\,dz=2\pi i\,\Res(f;z_0)=-\frac{2\pi i\,e^{i\pi/n}}{n}
\]
We can now consider the three pieces for $\alpha\to\infty$:

\begin{itemize}
    \item \underline{Outgoing segment:} $\displaystyle\int_{0}^{\alpha}\frac{dx}{x^n+1}\to I$, which is the desired integral.

    \item \underline{Arc:} Using Prop~3.26(a), since $n\ge 2$, we have,
    \[
    \left|\int_{C_\alpha}f\right|\le\frac{1}{\alpha^n-1}\cdot\frac{2\pi\alpha}{n}
    \xrightarrow[\alpha\to\infty]{}0
    \]

    \item \underline{Returning segmnet:} We can paramatrize $z=re^{2\pi i/n}$ with $r$ from $\alpha$ to
    $0$. We get $z^n=r^n e^{2\pi i}=r^n$ and $dz=e^{2\pi i/n}\,dr$, therefore,
    \[
    \int_{\alpha}^{0}\frac{e^{2\pi i/n}}{r^n+1}\,dr
    =-e^{2\pi i/n}\int_{0}^{\alpha}\frac{dr}{r^n+1}
    \xrightarrow[\alpha\to\infty]{}-e^{2\pi i/n}\,I
    \]
\end{itemize}

We can add the three contributions and we see that as $\alpha\to\infty$,
\[
I\bigl(1-e^{2\pi i/n}\bigr)=-\frac{2\pi i\,e^{i\pi/n}}{n}
\]
We can factor $1-e^{2\pi i/n}=-e^{i\pi/n}\bigl(e^{i\pi/n}-e^{-i\pi/n}\bigr)
=-e^{i\pi/n}\cdot 2i\sin(\pi/n)$. Therefore, we have,
\[
I=\frac{-\,2\pi i\,e^{i\pi/n}}{-\,2i\,n\,e^{i\pi/n}\sin(\pi/n)}
=\frac{\pi}{n\sin(\pi/n)}
\]
which finally gives,
\[
\int_{0}^{\infty}\frac{dx}{x^n+1}=\frac{\pi}{n\sin(\pi/n)}.
\]

\newpage
\section*{Ex 9.3.14}

\begin{prob}
Use Proposition 9.33 with $f(z)=(z^2-a^2)^{-1}$ to show that if $a\notin\Z$, then
\[
\pi\cot(\pi a)=\frac{1}{a}+\sum_{n=1}^{\infty}\frac{2a}{a^2-n^2}.
\]
\end{prob}

\textit{Pf.} Let $f(z)=\dfrac{1}{z^2-a^2}=\dfrac{p(z)}{q(z)}$ with $\deg p=0$ and
$\deg q=2$, so that $\deg p+2\le\deg q$. We have that its poles $\pm a$ are disjoint from
$\Z$ because $a\notin\Z$. Therefore, we can use Prop~9.33  to $g(z)=\pi f(z)\cot(\pi z)$ to obtain,
\[
\sum_{k=-\infty}^{\infty}\frac{1}{k^2-a^2}
=-\sum_{q(z)=0}\Res(g;z)=-\bigl(\Res(g;a)+\Res(g;-a)\bigr)
\]
Each of these poles is simple. Since $a\notin\Z$, $\cot(\pi z)$ is holomorphic and finite
at $\pm a$, so using Lemm~9.16 (Simple Pole Lemma) with denominator $z^2-a^2$ of derivative $2z$, we have,
\begin{align*}
\Res(g;a) &=\frac{\pi\cot(\pi a)}{2a} \\
\Res(g;-a) &=\frac{\pi\cot(-\pi a)}{-2a}=\frac{\pi\cot(\pi a)}{2a}
\end{align*}
using the fact that $\cot$ is odd. Therefore, we have,
\[
\sum_{k=-\infty}^{\infty}\frac{1}{k^2-a^2}=-\frac{\pi\cot(\pi a)}{a}
\]
On LHS, we can separate $k=0$ and pair $k$ with $-k$ (for $k\ge 1$), using $k^2=(-k)^2$, which gives,
\[
\sum_{k=-\infty}^{\infty}\frac{1}{k^2-a^2}
=-\frac{1}{a^2}+\sum_{n=1}^{\infty}\frac{2}{n^2-a^2}
\]
Therefore, $-\dfrac{1}{a^2}+\sum_{n\ge 1}\dfrac{2}{n^2-a^2}
=-\dfrac{\pi\cot(\pi a)}{a}$. We can multiply by $-a$ and use
$\dfrac{-2a}{n^2-a^2}=\dfrac{2a}{a^2-n^2}$ to get,
\[
\pi\cot(\pi a)=\frac{1}{a}+\sum_{n=1}^{\infty}\frac{2a}{a^2-n^2}
\]
as required. \qed

\newpage
\section*{Ex 11.1.3}

\begin{prob}
Complete the proof of Lemma 11.8.

\medskip
\emph{Lemma 11.8.} Suppose that $|z|<\tfrac12$. Then
\[
\tfrac12|z|\le|\Log(1+z)|\le\tfrac32|z|.
\]
\end{prob}

\textit{Pf.} From Maclaurin series $\Log(1+z)=\sum_{k\ge 1}(-1)^{k-1}z^k/k$, valid for $|z|<1$. We can divide by $z$ ($z\ne 0$) and compare coefficients with $\tfrac12$.
Comparing the bounds in equation~(11.2): for $0<|z|<1$, we have,
\[
\left|\frac{\Log(1+z)}{z}-1\right|
=\left|-\frac{z}{2}+\frac{z^2}{3}-\cdots\right|
\le\frac12\bigl(|z|+|z|^2+\cdots\bigr)=\frac12\frac{|z|}{1-|z|}
\]
We can complete the argument. Consider $0<|z|<\tfrac12$. The function
$t\mapsto t/(1-t)$ is increasing on $[0,1)$, so
\[
\frac{|z|}{1-|z|}<\frac{1/2}{1-1/2}=1
\]
when,
\[
\left|\frac{\Log(1+z)}{z}-1\right|\le\frac12
\]
Using the (reverse) triangle inequality, we have,
\[
1-\frac12\le\left|\frac{\Log(1+z)}{z}\right|\le 1+\frac12
\]
i.e.
\[
\frac12\le\left|\frac{\Log(1+z)}{z}\right|\le\frac32
\]
We can multiply through by $|z|>0$ to get,
\[
\tfrac12|z|\le|\Log(1+z)|\le\tfrac32|z|
\]
Using $z=0$, the three quantities equal $0$, so the inequalities hold there as well. This gives the assertion of Lemm~11.8.

\end{document}
